\subsection{分次交换群}

我们要把有关直和的定义（§1，第 8 号）翻译成另一种语言。

定义1. 给定一个写成加法形式的交换群 G 和一个集合 \(\Delta\) ， G 上一个类型为 \(\Delta\) 的分次是指子群的一个族 \((\mathbf{G}_{\lambda})_{\lambda \in L}\) ，G 是这些子群的直和。集合 G 连同由其群规律及其分次所定义的结构，称为一个类型为 \(\Delta\) 的分次（交换）群。

\(\Delta\) 称为 \(\mathbf{G}\) 的次数集合。一个元素 \(x \in \mathbf{G}\) 如果属于某个 \(\mathbf{G}_{\lambda}\) ，就称为齐次的；其次数为 \(\lambda\) ，倘若 \(x \in \mathbf{G}_{\lambda}\) 。因此元素 0 是一切次数的齐次元素；但若 \(x \neq 0\) 是齐次的，它只属于唯一的一个 \(\mathbf{G}_{\lambda}\) ；此时指标 \(\lambda\) （它使 \(x \in \mathbf{G}_{\lambda}\) ）称为 \(x\) 的次数（有时也称 \(x\) 的权），有时记作 \(\deg(x)\) 。每个 \(y \in \mathbf{G}\)

都可以唯一地写成齐次元素之和 \(\sum_{\lambda} y_{\lambda}\) ，其中 \(y_{\lambda} \in G_{\lambda}\) ； \(y_{\lambda}\) 称为次数为 \(\lambda\) 的齐次分量（或简称次数为 \(\lambda\) 的分量），是 \(y\) 的分量。当用“权”一词代替“次数”时，形容词“齐次的”换成“等权的”。

例. (1) 给定任意交换幺半群 \(\Delta\) （以 0 为单位元）和交换群 \(G\) ，分次 \((G_{\lambda})_{\lambda \in \Delta}\) 就定义在 \(G\) 上，办法是取 \(G_0 = G\) 以及 \(G_{\lambda} = \{0\}\) （对 \(\lambda \neq 0\) ）；这个分次称为平凡分次。

\begin{enumerate}
\def\labelenumi{(\arabic{enumi})}
\setcounter{enumi}{1}
\tightlist
\item
  设 \(\Delta, \Delta'\) 是两个集合， \(\rho\) 是 \(\Delta\) 到 \(\Delta'\) 的一个映射。设 \((\mathbf{G}_{\lambda})_{\lambda \in \Delta}\) 是一个类型为 \(\Delta\) 的分次（在交换群 \(\mathbf{G}\) 上的）；对 \(\mu \in \Delta'\) ，令 \(\mathbf{G}_{\mu}'\) 为诸 \(\mathbf{G}_{\lambda}\) 之和（ \(\rho(\lambda) = \mu\) ）；显然 \((\mathbf{G}_{\mu}')_{\mu \in \Delta'}\) 是一个类型为 \(\Delta'\) 的分次（在 \(\mathbf{G}\) 上的），称为由 \((\mathbf{G}_{\lambda})\) 借助映射 \(\rho\) 导出的分次。
\end{enumerate}

当 \(\Delta\) 是一个写成加法形式的交换群且 \(\rho\) 是映射 \(\lambda \mapsto -\lambda\) （ \(\Delta\) 到其自身的）时， \((\mathbf{G}_{\mu}^{\prime})\) 称为 \((\mathbf{G}_{\lambda})\) 的相反分次。

\begin{enumerate}
\def\labelenumi{(\arabic{enumi})}
\setcounter{enumi}{2}
\item
  若 \(\Delta = \Delta_1 \times \Delta_2\) 是两个集合的积，类型为 \(\Delta\) 的分次称为类型 \(\Delta_1, \Delta_2\) 的双分次。对所有 \(\lambda \in \Delta_1\) ，令 \(G_{\lambda}' = \bigoplus_{\mu \in \Delta_2} G_{\lambda\mu}\) ；对所有 \(\mu \in \Delta_2\) ，令 \(G_{\mu}' = \bigoplus_{\lambda \in \Delta_1} G_{\lambda\mu}\) ；显然 \((G_{\lambda}')_{\lambda \in \Delta_1}\) 是类型 \(\Delta_1\) 的一个分次，而 \((G_{\mu}'')_{\mu \in \Delta_2}\) 是类型 \(\Delta_2\) 的一个分次（在 \(G\) 上的）；这些分次称为由双分次 \((G_{\lambda\mu})\) 导出的部分分次。注意 \(G_{\lambda\mu} = G_{\lambda}' \cap G_{\mu}'\) ；反之，若 \((G_{\lambda}')_{\lambda \in \Delta_1}\) 和 \((G_{\mu}'')_{\mu \in \Delta_2}\) 是 \(G\) 上的两个分次，使得 \(G\) 是诸 \(G_{\lambda\mu} = G_{\lambda}' \cap G_{\mu}'\) 的直和，则这些子群构成一个类型 \(\Delta_1, \Delta_2\) 的双分次（在 \(G\) 上的）， \((G_{\lambda}')\) 和 \((G_{\mu}'')\) 是它的部分分次。我们把 \(\Delta\) 是有限个集合之积的情形的推广留给读者。
\item
  设 \(\Delta_0\) 是一个写成加法形式的交换幺半群，其单位元记作 0；设 I 是任意集合，且 \(\Delta_0^{(I)} = \Delta\) 表示积 \(\Delta_0^I\) 中由具有有限支撑的族 \((\lambda_t)_{t \in I}\) 组成的子幺半群。设 \(\rho: \Delta \to \Delta_0\) 是 \(\Delta\) 到 \(\Delta_0\) 上的、由 \(\rho((\lambda_t)) = \sum_{t \in I} \lambda_t\) 定义的满射（余对角的）同态。由每个类型 \(\Delta\) 的分次可以导出一个类型 \(\Delta_0\) 的分次，借助 \(\rho\) （例 2）；它称为与给定的类型 \(\Delta\) 的“多重分次”相伴的全分次。
\end{enumerate}

本号的定义和例子可立即推广到 G 是未必交换的群的情形；只需处处把直和的概念换成“限制和”的概念（§1，第 6 号，注）。注意在这种情形下 \(G_{\lambda}\) 是 G 的正规子群，且对 \(\lambda \neq \mu\) ， \(G_{\lambda}\) 的每个元素与 \(G_{\mu}\) 的每个元素可交换。

